\documentclass[12pt]{article}
\usepackage[utf8]{inputenc}
\usepackage[T1]{fontenc}
\usepackage[margin=1in]{geometry}
\usepackage{setspace}
\usepackage{xcolor}
\usepackage[colorlinks=true, urlcolor=blue!60!black, linkcolor=blue!60!black]{hyperref}
\setlength{\parindent}{0em}
\setlength{\parskip}{0.45em}
\renewcommand{\baselinestretch}{1.0}
\begin{document}
\begin{center}
{\Large \textit{Allen v.\ Wright}, 468 U.S.\ 737 (1984)\par}
\vspace{0.3em}
{\normalsize Case Brief --- One-Page Edition\par}
\end{center}

\textbf{Memory Jogger.} Black parents could not sue the IRS over lax enforcement against discriminatory private schools: no standing for abstract stigma or for harms not ``fairly traceable'' to the IRS (\href{https://supreme.justia.com/cases/federal/us/468/737/}{Justia}).

\textbf{Facts.} The IRS denied tax exemptions to discriminatory private schools but required only a certified nondiscriminatory policy. Parents of Black children in desegregating districts sued nationwide, alleging stigma and impairment of desegregated public education; none had applied to a private school.

\textbf{Procedural History.} The district court dismissed; the D.C.\ Circuit reversed and enjoined the IRS.

\textbf{Votes.} Reversed 5--3: O'Connor for the Court; Brennan and Stevens (with Blackmun) dissented; Marshall did not participate.

\textbf{Holding.} No Article~III standing. Stigmatic injury counts only for those personally denied equal treatment, and the desegregation injury depends on independent third-party choices, so it is not fairly traceable to the IRS.

\textbf{Reasoning.} Standing rests on separation of powers; recognizing it here would make courts ``continuing monitors'' of the Executive. Under \textit{Simon}, whether revoked exemptions would change schools' policies or enrollment was speculative. Brennan accused the Court of using ``standing to slam the courthouse door''; Stevens said causation followed from ``elementary economics.''

\textbf{Examples.} \emph{Same-Side:} States seeking more immigration arrests (\textit{United States v.\ Texas}). \emph{Opposite-Side:} a student personally rejected by an academy because of race. \emph{Fence-Sitter:} fuel producers hurt by an electric-vehicle mandate (\textit{Diamond Alternative Energy}).

\textbf{Critique.} Critics say the Court manipulated causation to dodge the merits and locked regulatory beneficiaries out of court.

\textbf{Key Quotation.} ``More important, the law of Art.\ III standing is built on a single basic idea --- the idea of separation of powers.'' --- Justice O'Connor, 468 U.S.\ at 752.

\end{document}
